\section{Local Functional Equation}

To treat the local field of number fields, we fix the notation for sake of convience.

\begin{definition}
    Let \(F\) be a field, then we call \(F\) is a local field if it is one of the following types:
    \begin{enumerate}
        \item Archimedean local field: \(\rr\) or \(\cc\); 
        \item Non-archimedean local field: a finite extension of \(\qq_p\)
    \end{enumerate} 
    Equivalently, a local field is a completion of a number field at a place, and it is non-archimedean if the place is finite.
\end{definition}

In the case of non-archimedean local field, we use the following notation:
\[
\begin{array}{rcl}
\mathcal O &:& \text{the valuation ring} \\[6pt]
\mathcal O^\times &:& \text{the unit group of } \mathcal O \\[6pt]
\mathfrak p &:& \text{the maximal ideal of } \mathcal O \\[6pt]
\varpi &:& \text{a uniformizer} \\[6pt]
\mathfrak{d} &:& \text{the different of  } F/\qq_p\\[6pt]


\end{array}
\]

\begin{definition}
    Let \(F\) be a local field, then \(f: F \to \cc\) is called a Schwartz-Bruhat function if it is one of the following cases:
    \begin{itemize}
        \item When \(F\) is archimedean, \(f\) is a Schwartz function, which is a smooth and rapidly deacrasing function of class \(C^{\infty}\).
        \item When \(F\) is non-archimedean, \(f\) is locally constant and compactly supported.
    \end{itemize}
    Let \(\mathcal{S}(F)\) denotes the \(\cc\)-vector space of all Schwartz-Bruhat functions on \(F\).
\end{definition}

\begin{definition}
    A LCA group is a locally compact hausdorff abelian group.
\end{definition}

\subsection{Additive Character}

Let \(F\) be a local field, then in each case we can know that \((F,+)\) is a locally compact abelian group, which induces a Haar measure \(\mu\) on \(F\), then for any non-zero element \(a \in F\), we can define a new measure on \(F\) by pull back:
\[\forall E \subset F \text{ measurable }, \quad  \mu_a(E) = \mu(aE)  \]
and we can conculde the following result:

\begin{lemma} \label{normalizedmodule}
    \(\mu_a\) is a Haar measure on \((F,+)\), and there exists a constant factor \(|a|\) such that \(\mu_a = |a| \mu\), such a factor is independent of the choice of \(\mu\) with
    \[|a| = \begin{cases}
        \text{the standard absolute value if } F = \rr \\
        \text{the square of standard absolute value if } F = \cc \\
        N(\pp)^{-v_{\pp}(a)} \text{ if } F \text{ is non-archimedean}
    \end{cases}\]
\end{lemma}

\begin{proof}
    The case of \(F =\rr\) is clear. If \(F = \cc\), we take \(a = u + iv\) and \(\mu\) be the lebesgue measure, let \(z = x+iy\) and then \(az = (ux - vy) + i(vx+uy)\), and we consider the integration in \(\rr^2\), then the jacobian of the transformation is 
    \[|\det J_a| = \begin{vmatrix}
        u & -v \\
        v & u
    \end{vmatrix} = u^2 + v^2 = \|a\|^2 \]
    so that shows \(d(az) = \|a\|^2 dz\). If \(F\) is non-archimedean, we take \(\mu\) be the Haar measure such that \(\mu(\oo) = 1\), and we set \(q = N(\pp) = \#(\oo/\pp)\), then
    \[1 = \mu(\oo) = \mu(\bigsqcup_{i=1}^q x_i + \pp) = \sum_{i=1}^q \mu(x_i + \pp) = q \mu(\pp) \]
    so it shows that \(\mu(\pp) = 1/q\), and same method shows that \(\mu(\pp^n) = 1/q^n\). So for any \(a \in F^\times\), we can write \(a = \varpi^nu\) with \(n = v_{\pp}(a)\) and \(u \in \oo^\times\), then \(a\oo = \varpi^n\oo = \pp^n\) and then 
    \[\mu(a\oo)/\mu(\oo) = 1/q^n = N(\pp)^{-{v_{\pp}(a)}}\]
\end{proof}

The constant factor can be extended to the whole field \(F\) by setting \(|0| = 0\), and we call it \textbf{(normalized) module} on \(F\). A obeservation here is that module is a power of standard absolute value on \(F\), so \(a \mapsto |a|\) defines a a continous map from \(F\) to \([0,\infty)\).

\begin{remark}
    More generally, for any LCA group \((G,+)\) with Haar measure \(\mu\), if \(\varphi\) is a continuous automorphism of \(G\), then \(\mu \circ \varphi\) is also a Haar measure with \(\mu\circ \varphi = \mathrm{mod}_G(\varphi) \cdot \mu\) for some constant factor only depending on \(\varphi\), such a factor is called the module of \(\varphi\) on \(G\). Let \(\mathrm{Aut}(G)\) denote the group of all continuous automorphisms of \(G\), then \(\varphi \mapsto \mathrm{mod}_G(\varphi)\) is actually a homomorphism from \(\mathrm{Aut}(G)\) to \(\rr_{>0}\). In the case of local field \(F\), \(F^\times\) can be embedded into \(\mathrm{Aut}(F)\) by \(a \mapsto m_a\) with \(m_a(x) = ax\), and the module \(\mathrm{mod}_F(m_a)\) is exactly \(|a|\) defined in above lemma.
\end{remark}



\begin{definition}
    Additive character of a local field \(F\) is a continous group homomorphism \(\psi: (F,+) \to \sph^1\), and \(\hat{F}\) is the set of all additive characters of \(F\), it is called the dual group of \(F\).
\end{definition}

The dual group is indeed a group under pointwise multiplication, for topology on \(\hat{F}\), we consider the compact-open topology, i.e. the topology generated by
\[[K,U]= \{\chi \in \hat{F} : \chi(K) \subset U\}\]
for every compact subset \(K\) of \(F\) and every open subset \(U\) of \(\sph^1\). It is a natural choice for functiona space, so dual group is ineed a topological group, a remarkable result here is that local field is self-dual.

\begin{theorem}
    Let \(\psi\) be a non-trival additive character of \(F\), then the map
    \[\Psi: F \to \hat{F}, \quad a \mapsto \psi_a :x \mapsto \psi(ax)\]
    is an isomorphism of LCA groups.
\end{theorem}

\begin{proof}
    The proof depends on a fact that all local fields are a metrizable space such that compact subset is bounded. It is not hard to check that \(\Psi\) is a injective homomorphism, so firstly we prove that \(\Psi\) is continous: Let \([K,U]\) be an open basis of \(\hat{F}\), then \[\Psi^{-1}([K,U]) = \{a \in F : \psi_a(K) \subset U\} = \{a \in F : aK \subset \psi^{-1}(U)\}\]
    notice that \( \psi^{-1}(U)\) is a open since \(\psi\) is continous, so by continous map \(m(x,y) = xy\) on \(F\times F\), for any \(y \in K\) there exists an neiborhood \(A_y\) of \(y\) and a neiborhood \(B_y\) of \(a\) such that \(B_yA_y \subset \psi^{-1}(U)\). Let \(y\) runs over all \(K\), and by compactness of \(K\), there exists a finite index set \(I\) such that \(K \subset \cup_{i\in I}A_i\), and we take \(V_a = \cap_{i \in I}B_i\) as a neiborhood of \(a\) such that \(V_a \subset \Psi^{-1}[K,U]\), it shows that \(a\) is an interior and then \(\Psi^{-1}[K,U]\) is open.

    Secondly, we prove that \(\Psi\) is an homemorphism to its image, it is enough to prove the image of open ball \(B_{\epsilon}\) centred at 0 with radius \(\epsilon\) is open in \(\hat{F}\). Non-trival character \(\psi\) implies an element \(b \in F \) such that \(\psi(b) \neq 1\), and we construct a neiborhood of trivial character by \([\overline{B_R(0)},W] \cap \Psi(F)\), where \(W\) is an neiborhood of \(0\) not cotaining \(\psi(b)\), and \(R \geq \|b\|/\epsilon\), then this neiborhood is contained in \(\Psi(B_{\epsilon})\). For other non-trival character in image, a translation of group allows us to conclude the result.

    Finally, \(\Psi(F)\) is complete since \(F\) is complete, and it is closed since \(\hat{F}\) is hausdorff, so it is enough to prove that \(\Psi(F)\) is dense in \(\hat{F}\), that will imply \(\Psi(F) = \hat{F}\). In fact, \( x \in \Psi(F)^\perp \) implies that \(\Psi(ax) = 1\) for any \(a \in F\), and it shows that the annihilator is zero, so \(\Psi(F)\) is dense in \(\hat{F}\).
\end{proof}

\begin{remark}
    A LCA group is not necessary self-dual, for example the dual group of \(\zz\) is \(\sph^1\), and the dual group of \(\sph^1\) is \(\zz\).
\end{remark}

By the theory of harmonic analysis on LCA, if LCA group \(G\) is self-dual, then it is nice to identify the fourier tranform of a function \(f\) on \(G\) as a function on \(G\) itself, so it motivates us to define the transform on local field as follows:

\begin{definition}
    Fix a non-trival additive character \(\psi\) and a Haar measure \(\mu\) on \(F\), the Fourier transform of type \((\psi, \mu)\) of a function \(f \in L^1(F)\) is defined as
    \[\hat{f}(y) = \int_F f(x) \psi(xy) \, d\mu(x)\]
\end{definition}

There is no confusion to identify \(y\) as \(\psi_y\) by isomorphism, and the norm of \(\psi\) is always 1, so \(|\hat{f}(y)| \leq \|f\|_{L^1} < \infty \) ensures that the transform is well-defined.

\begin{definition} \label{standardcharacter}
    The standard additive character of \(F\) is defined as following:
    \begin{itemize}
        \item If \(F = \rr\), let \(\psi(x) = e^{-2\pi i x}\);
        \item If \(F = \qq_p\), let \(\psi(x) = e^{2\pi i \{x\}}\), where \(\{x\}\) is the fractional part of \(x\) defined by
        \[x = \sum_{n \geq N} a_np^n \implies \{x\} = \sum_{n = N}^{-1} a_np^n\]
        where \(\{x\} = 0\) if \(N \geq 0\), i.e. \(x \in \zz_p\).
        \item Let \(F_0 = \rr \text{ or } \qq_p\), and \(F\) be a finite extension of \(F_0\), let the standard additive character of \(F\) be the compistion 
        \[F \xrightarrow{\mathrm{Tr}_{F/F_0}} F_0 \xrightarrow{\psi_0} \sph^1\]
        where \(\psi_0\) is the standard additive character of \(F_0\).
    \end{itemize}
\end{definition}

In particular, if \(F =\qq\), the standard character is \(\psi(z) = e^{-2 \pi i(z + \bar{z})}\). In the case of \(F =\qq_p\), we have \(\{x\} \in \zz[1/p]\) and \(x - \{x\} \in \zz_p\), the intersection \(\zz_p \cap \zz[1/p] = \zz\) implies that standard character \(\psi\) is a non-trival homomoprhism. In the aspect of continuity, we notice that \(\zz_p \subset \ker \psi\) as an open subset containing \(0\), which shows that \(\psi\) is continous at \(0\), then \(\psi\) is continous on \(\qq_p\) since it is a topological group. Hence \(\psi\) is exactly a non-trival additive character of \(\qq_p\), It is not clear than the case \(F =\rr\).

From now on, we will always use standard character in the Fourier transform, so the transform will depend only on the choice of Haar measure, so a normalization can be done to obtain the inversion formula.

\begin{proposition} \label{selfdualmeasure}
    Let \(F\) be a local field, then there exists a unique self-dual Haar measure \(\mu\), that means if \(f \in \mathcal{S}(F)\), then \(\hat{f} \in \mathcal{S}(F)\) and \(\widehat{\hat{f}}(x) = f(-x)\), explicitly the measure is given as following:
    \[dx = \begin{cases}
        \text{lebesgue measure of } \rr &\text{If } F = \rr \\
        2 \times \text{lebesgue measure of } \cc &\text{If } F = \cc \\
        \text{the Haar measure such that } \mu(\mathcal O) = N(\mathfrak{d})^{-\frac{1}{2}} &\text{If } F \text{ is non-archimedean}
    \end{cases}\]

\end{proposition}

\begin{proof}
    For any Haar measure \(\mu\), \(\widehat{\widehat{f}}(x) = cf(-x)\) holds for some constant \(c\), so it is enough to choose a test function to modify. If \(F = \rr\), we choose \(f(x) = e^{-\pi x^2}\) and calculate the transform under the lebesgue measure, than \(\hat{f} = f\); If \(F = \cc\) we choose \(f(x)\), and calculate the transform under \(c\cdot dl\) (c times the lebesgue measure), than \(\hat{f} = \frac{c}{2}f\) shows \(2\,dl\) is the self-dual measure; If \(F\) is non-archimdean we choose \(f = \mathbf{1}_{\mathcal O}\), then under some haar measure \(\mu\) on a \(\widehat{\widehat{\mathbf{1}_{\oo}}} = \mu(\oo)\mu(\oo^{\perp})\mathbf{1}_{\oo}\), and we assume difference \(\mathfrak{d} = \pp^{m}\) to calculate
    \[\mu(\oo^{\perp}) = \mu(\mathfrak{d}^{-1}) = \mu(\varpi^{-m}\oo) = |\varpi|^{-m} N(\oo) = N(\mathfrak{d})\mu(\oo)\]
    so normization \(\mu(\oo)\mu(\oo^{\perp}) = 1\) gives the Haar measure.
\end{proof}

\subsection{Multiplicative Character}
The module on local field \(F\) induced a continous homomorphism
\[F^\times \to \rr_{>0}, \quad x \mapsto |x|\]
the kernel of this homomorphism will be important in the following discussion, and we denote it by \(U\), and we can conclude
\[U = \begin{cases}
    \{\pm 1\} \quad &\text{If } F = \rr \\
    \sph^1 \quad &\text{If } F = \cc \\
    \mathcal O^\times \quad &\text{If } F \text{ is non-archimedean}
\end{cases}\]
The group \(U\) as a subspace of local field \(F\) is compact in each case, and also open if the \(F\) is non-archimedean.




\begin{definition}
    a (multiplicative) quasi-character is a continous group homomorphism \(\chi: F^\times \to \cc^\times\), and the set of all quasi-characcters is denoted by \(\mathcal{X}(F^\times)\).

    a quasi-character \(\chi\) is called a (unitary) character if \(\mathrm{im}(\chi) \subset \sph^1\); a quasi-character \(\chi\) is called unramified if \(\chi\) is trival on \(U\).
\end{definition}

The vocabulary here follows Tate's thesis, here are the examples in the case of archimedean local field, which is the classical result in harmonic analysis.

\begin{example} \label{archimedeanquasicharacter}
    Every quasi-character of \(\rr^\times\) is of the form
    \[\chi_{\epsilon,s}(t) = \mathrm{sign}(t)^{\epsilon} \cdot |t|^s, \quad \epsilon \in \{0,1\}, s \in \cc\]
    Every quasi-character of \(\cc^\times\) is of the form
    \[\chi_{n,s}(z) = (\frac{z}{\|z\|})^n \cdot |z|^s, \quad n \in \zz, s \in \cc\]
    where \(\|z\|\) is the standard absolute value on \(\cc\), and in both case, the character is unitary if and only if \(s\) is a purely imaginary number.
\end{example}

By the example, it is natural to separate the quasi-character into two parts, the first part is a unitary character on units, and the second part is a power of module, in fact it is unramified and we can generalize the result.

\begin{lemma} \label{unramifiedquasicharacter}
    A quasi-character \(\chi \in \mathcal{X}(F^\times)\) is unramified if and only if it is of the form \(\chi(x) = |x|^s\) for some \(s \in \cc\). Let \(\mathcal{U}(F^\times)\) denote the set of all unramified quasi-characters on \(F^\times\), then it forms a subgroup of \(\mathcal{X}(F^\times)\) with isomorphism
    \[\mathcal{U}(F^\times) \cong \begin{cases}
        \cc \quad &\text{If } F \text{ is archimedean} \\
        \cc/\frac{2\pi i}{\log N(\pp)}\zz \quad &\text{If } F \text{ is non-archimedean}
    \end{cases}\]


\end{lemma}

\begin{proof}
    Notice that \(\chi\) is unramified if and only if \(\chi\) factors through \(|F^\times| \), because \(U \subset \ker \chi\) if \(\chi\) is trivial on \(U\), while \(U\) is the kernel of the surjection \(x \mapsto |x|\),  that means there exists a continous homomorphism \(\tilde{\chi}: |F^\times| \to \cc^\times\) such that the following diagram commutes: % https://q.uiver.app/#q=WzAsMyxbMCwwLCJGXlxcdGltZXMiXSxbMSwwLCJcXG1hdGhiYntDfV5cXHRpbWVzIl0sWzAsMSwifEZeXFx0aW1lc3wiXSxbMCwxLCJcXGNoaSJdLFswLDIsInhcXG1hcHN0byB8eHwiLDJdLFsyLDEsIlxcdGlsZGV7XFxjaGl9IiwyXV0=
\[\begin{tikzcd}
	{F^\times} & {\mathbb{C}^\times} \\
	{|F^\times|}
	\arrow["\chi", from=1-1, to=1-2]
	\arrow["{x\mapsto |x|}"', from=1-1, to=2-1]
	\arrow["{\tilde{\chi}}"', from=2-1, to=1-2]
\end{tikzcd}\]
    Hence if \(\chi\) is unramified, it is enough to prove that continous homomorphism \(\tilde{\chi}:|F^\times| \to \cc^\times\) is of the form \(x \mapsto x^s\) for some \(s \in \cc\). If \(F \) is archimedean, it is immediate by the example above, if \(F\) is non-archimedean, \(|F^\times| = q^{\mathbb{Z}}\) with \(q = N(\pp)\), so the morphism is determined by the image of \(q\), if we set \(\tilde{\chi}(q) = a\), then we can choose \(s = \log a \) with respect to a branch of logarithm. 

    Conversely, it is not hard to check that \(\chi_s(x) = |x|^s\) is unramified for any \(s \in \cc\), which proves the first part of the lemma and induces a surjective map
    \[\Phi: \cc \to \mathcal{U}(F^\times), \quad s \mapsto \chi_s\]
    and it is not hard to check \(\mathcal{U}(F^\times)\) is indeed a subgroup of \(\mathcal{X}(F^\times)\) and \(\chi_{s+z} = \chi_{s}\cdot\chi_{z}\) for any \(s,z \in \cc\), which implies that \(\Phi\) is a group homomorphism with kernel \(\{s \in \cc : |x|^s = 1 , \forall x \in F^\times \}\). If \(F\) is archimedean, the kernel is trivial so \(\Phi\) is an isomorphism; if \(F\) is non-archimedean, the kernel is exactly \(\{s \in \cc : |\varpi|^s = 1 \} = \frac{2\pi i}{\log q}\zz\), so it induces a desired isomorphism.
\end{proof}

\begin{lemma}
    The map \(\chi \mapsto (\chi|_U,{\chi} \cdot ({\chi|_U \circ \rho})^{-1}) \) gives an isomorphism
        \[\mathcal{X}(F^\times) \xrightarrow{\sim} \Hom_{\text{cont}}(U,\sph^1) \times \mathcal{U}(F^\times    )\]
    where \(\rho : F^\times \to U\) is a retraction depending on the local field \(F\), and every quasi-character \(\chi\) is of the form \(\chi = \tilde{\chi} |\cdot|^s\) for some \(s \in \cc\) and a unitary character \(\tilde{\chi}\) conciding with \(\chi\) on \(U\).
\end{lemma}

\begin{proof}
    It is enough to prove the non-arcimedean case, since we have seen that \(\rho(x) = x/\|x\|\) when \(F\) is archimedean. Notice that \(U = \mathcal{O}^\times\) is a compact subgroup of \(F^\times\), then the restriction \(\chi|_U\) must be a unitary, and we define \(\rho(x) = \varpi^{-v_{\pp}(x)}x\), which is a continous homomorphism with \(\rho(\varpi^nu) = u\) for any \(n \in \zz\) and \(u \in \oo^\times\). Hence \(c(x) = \chi \cdot (\chi|_U \circ \rho)^{-1}\) defines a quasi-character with \(c(u) = 1\) for any \(u \in \oo^\times\), by lemma \ref{unramifiedquasicharacter} it shows that \(c\) is unramified of the form \(c(x) = |x|^s\) for some \(s \in \cc\), so \(\chi = \chi_u  |\cdot|^s\) with \(\tilde{\chi} = \chi|_U \circ \rho\).

    Conversely, we can construct a inverse map by \((c,\chi_s) \mapsto (c \circ \rho)\cdot \chi_s\), which proves the map is an isomorphism. 
\end{proof}

\begin{remark}
    The isomorphism is also topological, the retraction \(\rho\) induces a map
    \[\rho^*: \Hom_{\text{cont}}(U,\cc^\times) \to \Hom_{\text{cont}}(F^\times,\cc^\times), \quad \chi \mapsto \chi \circ \rho\]
    such a map is continous since local field is locally compact and Hausdorff, in this case the map defined in the lemma is a homeomorphism.
\end{remark}



\begin{lemma}
    Let \(\chi:F^\times \to \cc^\times\) be a quasi-character of a non-archimedean local field \(F\), then there exists a integer \(n \geq 0\) such that \(\chi\) is trivial on \(1 + \pp^n\).
\end{lemma}
\begin{proof}
    \(\sph^1\) is a group has NSS (no small subgroup) properties, then there exists a neighborhood \(V\) of \(1\) in \(\sph^1\) such that \(V\) does not contain any non-trivial subgroup, then \(\chi^{-1}(V)\) is a open neighborhood of \(1\) in \(F^\times\). In fact, all open subgroups of \(\oo^\times\) form a chain  as following:
\[\oo^\times \varsupsetneqq 1+\pp \varsupsetneqq 1+\pp^2 \varsupsetneqq \cdots\]
and form a basis of open neighborhood of \(1\) in \(F^\times\), so there exists a integer \(n \geq 0\) such that \(1 + \pp^n \subset \chi^{-1}(V)\), which implies that \(\chi\) is trivial on \(1 + \pp^n\).
\end{proof}

Hence there is no confusion to take the samllest integer \(\mathfrak{f}\) and define it as the \textbf{conductor} of quasi-character \(\chi\), with a convention that \(1+\pp^0 = \oo^\times\), then \(\chi\) is unramified if and only if its conductor is \(0\), and such a concept measures how far a quasi-character is away from being unramified.

\begin{remark}
    Let \(F = \qq_p\) and \(\chi: F^\times \to \cc^\times\) be ramified character, and then \(\chi|_{\zz_p^\times}\) factors through \(\zz_p^\times/(1+p^\mathfrak{f}\zz_p)\) with conductor \(\mathfrak{f} > 0\) as following

    % https://q.uiver.app/#q=WzAsNCxbMSwwLCJcXFpfcF5cXHRpbWVzLygxK3BeblxcWl9wKSJdLFswLDAsIlxcWl9wIl0sWzIsMCwiKFxcWi9wXm4pXlxcdGltZXMiXSxbMywwLCJcXG1hdGhiYntTfV4xIl0sWzEsMF0sWzAsMiwiXFxzaW0iXSxbMiwzLCJcXHZhcnBoaSJdXQ==
\[\begin{tikzcd}
	{\zz_p^\times} & {\zz_p^\times/(1+p^\mathfrak{f}\zz_p)} & {(\zz/p^n)^\times} & {\mathbb{S}^1}
	\arrow[from=1-1, to=1-2]
	\arrow["\sim", from=1-2, to=1-3]
	\arrow["\varphi", from=1-3, to=1-4]
\end{tikzcd}\]
    the proof is omitted here, and the interesting result is that the ramified part of \(\chi\) is related to a primitive dirichlet character \(\varphi\) modulo \(p^{\mathfrak{f}}\). 
\end{remark}





Later we will see that the local zeta integral will be defined on the character group, so it is better to treat it as a complex manifold such that we can talk about the meromorphic continuation and functional equation.

\begin{definition}
    Let \(\chi\) and \(\chi'\) be two quasi-character of \(F^\times\), they are called equivalent if \(\chi'|_U = \chi'|_U\), or equivalently, \(\chi/\chi'\) is unramified.
\end{definition}

    It defines a equivalence relation on \(\mathcal{X}(F^\times)\), and the equivalence class of a quasi-character is \([\chi] = \{\chi|\cdot|^s : s \in \cc\} = \chi \cdot \mathcal{U}(F^\times)\), and lemma \ref{unramifiedquasicharacter} shows that \(\mathcal{U}(F^\times)\) is indeed a complex lie group of dimension 1, so each class can be identified as a Riemann surface by translation of topological group, and thus \(\mathcal{X}(F^\times)\) can be identified as a disjoint union of Riemann surfaces.

Finally, to consider the continuation of local zeta function, Tate introduces a notation to denote the "dual" of a quasi-character as following:

\begin{definition}
    Let \(\chi \in \mathcal{X}(F^\times)\), its \textbf{twisted dual} is defined as \(\chi^\vee = \chi^{-1}|\cdot|\).
\end{definition}

It is not hard to check that \(\chi \mapsto \chi^\vee\) is a homeomorphism on \(\mathcal{X}(F^\times)\) with \((\chi^\vee)^\vee = \chi\), and the image of a connected component \([\chi]\) is exactly the component \([\chi^{-1}]\).














\subsection{Local Zeta Function}
 Now we will consider the integrand on \(F^\times\), which will be used to define the local zeta integral.

\begin{lemma} \label{local-multiplicative-haarmeasure}
    Let \(dx\) be the self-dual Haar measure on \(F\) defined in the previous, then 
    \begin{enumerate} [(a)]
        \item \(m(E) = \int_{E}\frac{dx}{|x|}\) defines a Haar measure on multiplicative group \(F^\times\), and we denote it by \(d^\times x\);
        
        \item If\(F\) is non-archimedean, \(\oo^\times \) gets measure \((1-N(\pp)^{-1})N(\mathfrak{d})^{-1/2}\).
        \item A function \(f\) is integrable on \(F^\times\) with respect to \(d^\times x\) if and only if \(x \mapsto f(x)/|x|\) is integrable on \(F-{0}\) with respect to \(dx\). 
    \end{enumerate}
\end{lemma}
\begin{proof}
    
\end{proof}


\begin{definition}
    Fix a function \(f \in \mathcal{S}(F)\), the local zeta function associated with \(f\) is defined as a integrand on space \(\mathcal{X}(F^\times)\) by
\[Z(f,\chi) = \int_{F^\times} f(x)\chi(x) \, d^\times x\]
Let \(\chi\) be a quasi-character of the class \([\chi_0]\) of the form \(\chi_0|\cdot|^s\), it is equivalent to write
\[Z(f,\chi_0,s) = \int_{F^\times} f(x)\chi_0(x)|x|^s d^\times x\]
and we call here \(s\) as the complex parameter of \(\chi\) and \(\sigma = \mathrm{Re}(s)\) as the \textbf{exponent} of \(\chi\).
\end{definition}

Review that a function on Riemann surface \(f: X \to \cc\) is called holomorphic (resp. meomorphically) at \(s \in X\) if there exists a local chart \((U,\phi)\) with \(s \in U\) such that \(f \circ \phi^{-1}\) is holomorphi (resp. meomorphically) at \(\phi(s)\) as a complex function.

\begin{lemma} \label{holomorphic-halfplane}
    Following the notation above, \(\chi \mapsto Z(f,\chi)\) defines a holomorphic function on the domain of quasi-character with exponent \(\sigma > 0 \).
\end{lemma}

\begin{proof}
    We assume that \(K\) is a compact subset of the domain, for the archimedean case, schwartz condition implies that \(f\) is bounded by a constant \(B\) and \(f\) rapidly decreasing faster than any power of \(1/|x|\) at infinity, so an estimate can be done by
    \[|Z(f,\chi)| \leq B \int_{0< |x| \leq 1}|x|^{\sigma -1} \,d x + C_N\int_{|x| \geq 1} \frac{|x|^{\sigma-1} d x}{(1+ |x|)^N}\]
    \(\sigma>0\) ensures the first integral is finite, and the compactness of \(K\) ensures the existence of \(N\) such that \(N > \sigma\) and the second integral is finite. If \(F\) is non-archimdean, suppose that \(\mathrm{supp}(f) \subset \varpi^N \oo\) for some integer \(N \geq 0\), and then 
    \[|Z(f,\chi)| \leq C \int_{0<|x|<N(\pp)^N}|x|^{\sigma} d^\times x  = C \sum_{n \geq -N}\int_{\varpi^n \oo^\times}|x|^{\sigma} d^\times x\]
    and lemma \ref{local-multiplicative-haarmeasure} (a) allows us to set \(\mathrm{Vol}(\oo^\times) = (1-N(\pp)^{-1})N(\mathfrak{d})^{-1/2}\), then
    \[|Z(f,\chi)| \leq C \cdot \mathrm{Vol}(\oo^\times) \sum_{n \geq -N}N(\pp)^{-n \sigma}  = C \cdot \mathrm{Vol}(\oo^\times) \frac{N(\pp)^{N\sigma}}{1-N(\pp)^{-\sigma}}\]
    the last equality holds when \(\sigma>0\) by geometric series. Finally, the inequalities above is enough to show the uniform convergence of \(\chi \mapsto Z(f,\chi)\) in each compact of the component with \(\sigma>0\), so in each component the function is holomorphic on the domain  \(\sigma>0\).
\end{proof}

\begin{lemma}[Tate] \label{tate}
    Fix any two functions \(f,g \in \mathcal{S}(F)\), and let \(\chi\) be a quasi-character \(\chi\) of exponent \(0< \sigma <1\), then \(\chi^\vee\) is of exponent \(0 < \sigma <1\) as well, and the following identity holds
    \[Z(f,\chi)Z(\hat{g},\chi^\vee) = Z(\hat{f},\chi^\vee)Z(g,\chi)\]
\end{lemma}

\begin{proof}
    Applying the Fubini's theorem to separate the integral 
    \[Z(f,\chi)Z(\hat{g},\chi^\vee) = \int_{(F^\times)}f(x)g(z)\chi(xy^{-1})\psi(yz)|yz| \, d^\times x \,d^\times y \, d^\times z\]
    and then we change variable \(y =xt\) and then \(u = zt\) to calculate
    \begin{align*}
        Z(f,\chi)Z(\hat{g},\chi^\vee) &= \int_{(F^\times)^3} f(x)g(z)\chi(t^{-1})\psi(txz)|txz| \, d^\times x \,d^\times t \, d^\times z \\
        &= (\int_{F^\times}g(z)\chi(z) \, d^\times z )(\int_{(F^\times)^2} f(x)\chi(u)^{-1}\psi(ux)|ux| \, d^\times x \,d^\times u) \\
        &= Z(g,\chi)Z(\hat{f},\chi^\vee)
    \end{align*}
    notice that in finally equality we use the definition \(|ux|d^\times x = |u|dx\) to recover the local zeta function.
\end{proof}

Loosely speaking, the last lemma means the quotient of \(Z(f,\chi)\) and \(Z(\hat{f},\chi^\vee)\) is independent of the choice of function, so we can definie a meromorphic function on domain of quasi-character of exponent \(0 < \sigma <1\) by
\[\rho(\chi) = \frac{Z(\hat{f},\chi^\vee)}{Z(f,\chi)}\]
the choice of \(f\) should satisfy \(\chi \mapsto Z(f,\chi) \) is not zero function, and we call it \(\rho\)-factors of local field \(F\). The analytic continuation of \(\rho\)-factors will be the key to prove the local functional equation, Tate's strategy is to calculate each case explicitly.

\begin{example}
    If \(F =\rr\), by example \ref{archimedeanquasicharacter} we have
    \[\mathcal{X}(\rr^\times) = [\chi_0] \sqcup [\chi_1]\]
    with \(\chi_0\) the trival character and \(\chi_1\) the sign character, and they are representative of even characters and odd characters respectively. In the component \([\chi_0]\), we can choose \(f_0(x) = e^{-\pi x^2}\) with \(\hat{f_0} = f_0\); In the compoenent \([\chi_1]\), the choice of \(f_0\) is not proper since its local zeta function will vanish, so we modify it by \(f_1(x) = x e^{-\pi x^2}\) with \(\hat{f_1} = i\hat{f_1}\). After some calculation we can conclude the local zeta function
    \[Z(f_{\epsilon},\chi_{\epsilon,s}) = \pi^{-\frac{s+\epsilon}{2}}\Gamma\!\left(\frac{s+\epsilon}{2}\right)\]
    and its \(\rho\)-factor
    \[\rho(\chi_{\epsilon,s}) = i^{\epsilon} 2^{1-s} \pi^s \cos\frac{\pi(s+\epsilon)}{2} \, \Gamma(s)\]
    for quasi-character \(\chi_{\epsilon,s}(t) = \mathrm{sign}(t)^{\epsilon}|\cdot|^s\) of exponent \(\sigma = \mathrm{Re}(s) \in (0,1)\).
\end{example}


\begin{example}
    If \(F = \cc\), by example \ref{archimedeanquasicharacter} we have
    \[\mathcal{X}(\cc^\times) = \bigsqcup_{n \in \zz} [c_n]\]
    with \(c_n(z) = (z/\|z\|)^n\) the unitary character, and we can choose schwartz functions \(f_n\) for each component \([c_n]\) by 
    \[f_n(z) = \begin{cases}
        (\bar{z})^n e^{-2\pi \|z\|^2} & n \geq 0 \\
        z^{-n} e^{-2\pi \|z\|^2} & n < 0
    \end{cases}\]
    with fourier tranform \(\hat{f_n} = i^{|n|}f_{-n}\), and in each case the local zeta function does not vanish, which allows us to calculate the local zeta function
    \[Z(f_n,\chi_{n,s}) = (2\pi)^{1-s+\frac{|n|}{2}} \, \Gamma(s+ \frac{|n|}{2})\]
    and its \(\rho\)-factor
    \[\rho(\chi_{n,s}) = (-i)^{|n|}(2\pi)^{1-2s} \frac{\Gamma(s+\frac{|n|}{2})}{\Gamma(1-s+\frac{|n|}{2})}\]
    for quasi-character \(\chi_{n,s} = c_n |\cdot|^s\) of exponent \(\sigma \in (0,1)\).
\end{example}

\begin{example}
    If \(F\) is non-archimedean, let \(\mathcal{X}_n(F^\times)\) denote the set of quasi-characters of conductor \(n\), then we have a decomposition
    \[\mathcal{X}(F^\times) = \bigsqcup_{n \geq 0} \mathcal{X}_n(F^\times)\]
    where \(\mathcal{X}_0(F^\times)\) is just the set of unramified quasi-characters. For each class \(\mathcal{X}_n(F^\times)\), we can choose a test function \(f_n \in \mathcal{S}(F)\) by
    \[f_n(x) = \begin{cases}
        \psi(x) = \exp (2 \pi i \, \{\mathrm{Tr}_{F/\qq_p}(x)\}) & \text{if } x \in \mathfrak{d}^{-1}\pp^{-n} \\
        0 & \text{otherwise}
    \end{cases}\]
    and its transform are 
        \[\hat{f}_n(x) = \begin{cases}
        (N(\mathfrak{d})^{1/2} N(\pp)^n)  & \text{if } x \in 1+ \pp^n \\
        0 & \text{otherwise}
    \end{cases}\]
\end{example}






\begin{corollary}
    The \(\rho\)-factor of the unramified quasi-character is 
    \[\rho(|\cdot|^s) = N(\mathfrak{d})^{s-1/2}\frac{1-N(\pp)^{s-1}}{1-N(\pp^{-s})}\]
    If \(c_n\) is a unitary character of conductor \(n > 0\), the \(\rho\)-factor of the ramified quasi-character is 
    \[\rho(c_n|\cdot|^s) = N(\mathfrak{d} \pp^n)^{s-\frac{1}{2}} W(\chi)\]
    with root number \(W(\chi)\) defined by
    \[W(\chi) = N(\pp)^{-n/2} \sum_{\bar{a} \in \oo^\times/1+\pp^n } \chi(a) \psi(\frac{a}{\varpi^{d+n}})\]
\end{corollary}












\begin{theorem}[Tate's local functional equation] $ \\$
    Let \(F\) be a local field and fix a function \(f \in \mathcal{S}(F)\), then \(\chi \mapsto Z(f,\chi)\) has a meromorphic continuation to whole space  \(\mathcal{X}(F^\times)\), and there exists a meromorphic function \(\rho: \mathcal{X}(F^\times) \to \cc\) such that for any quasi-character \(\chi\) such that , the following functional equation holds
    \[Z(\hat{f},\chi^\vee) =\rho(\chi) Z(f,\chi) \]
    where \(\rho\) is independent of the choice of \(f\) and 
\end{theorem}

\begin{proof}
    In each case of local field, \(\rho(\chi) = Z(\hat{f_C},\chi^{\vee})/Z(f_C,\chi)\) with \(f_C\) choosen as above examples, so \(\rho\) has a meomorphic continuation to whole space \(\mathcal{X}(F^\times)\). Let \(h(\chi) = \rho(\chi)Z(f,\chi)\), lemma \ref{holomorphic-halfplane} shows that \(h\) is meomorphic on the domain of exponent \(\sigma > 0\), and by lemma \ref{tate} we can calculate that \(h(\chi) = Z(\hat{f},\chi^\vee)\) in the domain of exponent \(0 < \sigma < 1\), and \(\chi \mapsto Z(\hat{f},\chi^\vee)\) defines a holomorphic function on the domain of exponent \(\sigma < 1\). Hence the continuation of \(\chi \mapsto Z(f,\chi)\) can be concluded by the continuation of \(h\) by the uniqeness of meromorphic continuation.
\end{proof}


\begin{example}
    Some calculation can be done for some sepecific example, we can consider a number field, its completion are \(r_1\) real places and \(r_2\) complex places, and the rest are the local field at some prime ideaal \(\pp\).
    If \(F = \qq\) and we choose \(f(x) = e^{-\pi x^2}\), then we can calculate
    \[Z(f,s) = 2\int_{0}^\infty e^{\pi x^2}x^{s-1} \, dx = \pi^{-s/2}\Gamma(s/2)\]

    If \(F = \qq_p\) and we choose \(f(x) = \mathbf{1}_{\zz_p}\), then use the fact that \(\zz_p-\{0\} = \bigcup_{n\geq 0 } p^n\zz_p^\times\) we can calculate
    \begin{align*}
        Z(f,s) &= \sum_{n \geq 0} \int_{p^n\zz_p^\times} |x|_p^s \, d^\times x \\
        &= \sum_{n \geq 0} p^{-ns} \int_{\zz_p^\times} \, d^\times x \\
        &= \sum_{n \geq 0} p^{-ns} = \frac{1}{1-p^{-s}}
    \end{align*}
    Pay attention that in both case, the last equality holds when \(\mathrm{Re}(s)> 0 \), and by Euler's product formula we can find that the proudct of all local zeta functions is so-called completed zeta function
    \[s \mapsto \pi^{-s/2}\Gamma(s/2) \zeta(s)\]
    which plays an important role in the functional equation of Riemann zeta function, it can go back to Riemann's famous short paper ~\cite{Ri}.
\end{example}

In the example above, variable \(s\) can be seen as a parameter of unramified quasi-character \(\chi_s = |\cdot|^s\), and beyond the example, we should consider other type of zeta function. For instance, a dirichlet \(L\)-function asspciated with a dirichlet character \(\chi\), which is defined by the following Euler product
\[L(s,\chi) = \prod_{p \text{ prime }} \frac{1}{1-\chi(p)p^{-s}}\]
it suggests that we should consider all type of quasi-character.




